\section{SMART goal 2: Tool decision}
As part of SMART goal 2, the DQM team was followed on their way towards a tool decision for \gls{DQM}. The work happened in paralell to the first goal. For the project team, this goal serves to find the tool used in the follow up of the master thesis and allows to evaluate its possibilities right away within SMART goal 3.


\subsection*{Tool evaluation}
Katze-.---

\subsection*{Tool selection for DWH}
The DWH Decided that Cats are to hangry beasts

\subsection*{Tool selection for Swiss Post}
The decision for Swiss Post was taken..
Not the dqm plattform..


\subsection*{DQM plattform tool}
The decision for Swiss Post was taken..
Not the dqm plattform..
18.05.2026  | Meeting


To-Dos
Follow-up Tasks (wichtigster Teil)
1. Rule-Based DQM testen
Injection-/Review-Tests durchführen bzw. Resultate prüfen.
Validieren, dass die DQM-Regeln korrekt funktionieren.
2. Nach Deployment Status prüfen
Push/Deployment durchführen.
Status der Glue Jobs kontrollieren.
Fehler und Laufzeiten überprüfen.
3. Power BI Integration klären
Prüfen, wie die verfügbaren Daten aus SyncMetrics in Power BI eingebunden werden können.
Verantwortliche Power-BI-Entwickler identifizieren/einbeziehen.
4. Ownership festlegen
Verantwortlichkeit für die Power-BI-Implementierung definieren.
Klären, wer die nächsten Umsetzungsschritte startet.
5. Datenzugang prüfen
Verifizieren, warum die Daten nur lokal verfügbar sind.
Prüfen, ob ein zentraler bzw. servicebasierter Zugriff (anstelle lokaler Nutzung) möglich ist.


\subsection*{Conclusion}





Update on 01.06.2026
The weekly coordnation calls have been calceled

While the rule based DQM is to be used for AWS,
The team searching company wide DQM System, continues looking for another solution

Thus coordination calls no longer needed between the two teams

\subsection*{Discussion with Solution Architect}
The decision for Swiss Post was taken..
Not the dqm plattform..

In call with architect, i got the hint to use the Monitoring built by AWS Modernisation

Call mit Solution      Architect, Abklärung mit welchen Tools für das Monitoring gearbeitet wird:
Ziel davon : Das       richtige Tool verwenden
Oder Abklären ob eine       andere Idee im Raum steht
Verwendung des       Monitoring Setups möglich
Monitoring Setup wurde        im Rahmen von dem Projekt «AWS Modernisation» bei der Post aufgebaut
Monitoring sowohl für       Job Monitoring als auch Datenverlauf vorgesehen
Möglichkeit als PoC für       Datenmonitoring zu prüfen wie/ob man Daten aus der Sendungsverarbeitung       in das Monitoring bekommt
Fazit       : Die       Aufgabe könnte für eine Master Arbeit taugen. Und nach einem PoC weitere       Schritte bis integration des Datenmonitoring und ausweitung davon führen.